% Generated draft: requires mathematical and source review.
\documentclass[11pt]{article}
\usepackage[margin=25mm]{geometry}
\usepackage{fontspec}
\setmainfont{Latin Modern Roman}
\usepackage{amsmath,amssymb,mathtools,unicode-math}
\setmathfont{Latin Modern Math}
\usepackage{enumitem,needspace,xurl,hyperref}
\setlength{\parindent}{0pt}
\setlength{\parskip}{5pt}
\setlength{\emergencystretch}{3em}
\newcommand{\N}{\mathbb{N}}
\newcommand{\Z}{\mathbb{Z}}
\newcommand{\Q}{\mathbb{Q}}
\newcommand{\R}{\mathbb{R}}
\newcommand{\C}{\mathbb{C}}
\newcommand{\F}{\mathbb{F}}
\newcommand{\A}{\mathbb{A}}
\newcommand{\PP}{\mathbb{P}}
\newcommand{\E}{\mathbb{E}}
\newcommand{\eps}{\varepsilon}
\newcommand{\dd}{\,\mathrm{d}}
\begin{document}
\section*{UnsolvedMath ID: 800204}
\textbf{Source problem code:} catalogue entry
\subsection{Definitions and mathematical statement}
A holomorphic contact structure on a smooth complex manifold \(X\) of dimension \(2n+1\) is a holomorphic subbundle \(\mathcal H\subset T_X\) of rank \(2n\), with quotient line bundle \(L=T_X/\mathcal H\), such that the bracket pairing
\[
\mathcal H\otimes\mathcal H\longrightarrow L,\qquad (u,v)\longmapsto [u,v]\bmod\mathcal H
\]
is nondegenerate at every point. A projective contact manifold is a smooth connected complex projective manifold equipped with such a structure.

For a compact complex manifold \(M\), let \(\mathbb P(T_M^*)\) denote the bundle of one-dimensional linear subspaces of its cotangent spaces. It carries the natural contact distribution whose fiber at \((m,[\xi])\) consists of those tangent vectors \(v\) satisfying \(\xi(\dd\pi(v))=0\), where \(\pi:\mathbb P(T_M^*)\to M\) is the projection.

For a finite-dimensional nonabelian simple complex Lie algebra \(\mathfrak g\), let \(G\) be its adjoint group and let \(X_{\mathfrak g}\subset\mathbb P(\mathfrak g)\) be its unique closed \(G\)-orbit. As described in the supplied source, every other orbit contains \(X_{\mathfrak g}\) in its closure, and \(X_{\mathfrak g}\) is a contact Fano manifold; here Fano means that the anticanonical line bundle is ample.

\textbf{Conjecture 13 (folklore, as listed by Beauville in 2010).} For every projective contact manifold \(X\), either
\[
X\cong\mathbb P(T_M^*)\quad\text{for some compact complex manifold }M,
\]
or
\[
X\cong X_{\mathfrak g}\quad\text{for some simple complex Lie algebra }\mathfrak g.
\]
The first alternative retains the source's hypothesis that \(M\) is compact complex; no additional projectivity hypothesis on \(M\) is imposed here.

The displayed isomorphisms express a classification of the underlying complex manifolds. The supplied wording does not explicitly say whether they must also identify the given contact distributions. It also gives no lower dimension bound. The contact definition above can be used for \(n\geq0\); if contact manifolds are conventionally required to have dimension at least three, restrict to \(n\geq1\). These interpretive qualifications require review.
\subsection{Short English statement}
The folklore conjecture listed by Beauville says that every projective holomorphic contact manifold is a projectivized cotangent bundle of a compact complex manifold or the closed adjoint orbit of a simple complex Lie algebra. The intended dimension convention and whether the classification preserves contact structures are not explicit in the supplied record.
\subsection{Sources}
\begin{itemize}
\item Local UnsolvedMath registry, permanent ID \(800204\), exact source problem code \texttt{catalogue entry}, title ``Holomorphic Symplectic Geometry --- Conjecture 13.'' The supplied registry labels the problem \texttt{open}; this label is not an independently verified mathematical or literature assessment.
\item Beauville, \emph{Holomorphic Symplectic Geometry: A Problem List} (2010), Conjecture 13. Supplied source URL: \url{https://arxiv.org/abs/1002.4321}. The record describes the conjecture as folklore and lists Beauville as proposer, with proposed year 2010.
\item Supplied provenance: extraction method \texttt{arxiv\_latex}; recorded access date 2026-10-16. This is source-provided provenance, not a claim of independent access. The introductory assertion that only the two displayed classes were known is retained as historical context from the supplied source, not as a current literature assessment.
\end{itemize}
\end{document}
